% Part: computability
% Chapter: recursive-functions
% Section: partial-functions

\documentclass[../../../include/open-logic-section]{subfiles}

\begin{document}

\olfileid{cmp}{rec}{par}
\olsection{الدوال العودية الجزئية}

قبل تعريف الدوال العودية، لننظر فيما زاد برهاننا على مجرد إثبات دوال قابلة
للحساب غير عودية بدائية. فأصل الحجة
إمكان تعداد دوال $f_0,f_1,\dots$ على نحو يجعل،
بالنظر إلى الحجتين $x$ و$y$، قيمة $f_x(y)$ قابلة للحساب. فالحجة إذن جارية في
\emph{كل صنف من الدوال يمكن تعداده على هذا النحو}. وهنا
موضع الإشكال: نطلب وصفًا صريحًا للدوال القابلة للحساب، ثم لا يكون الوصف
الذي يتيح هذا التعداد الفعال لدوال كلية مستنفدًا لها!

ويدفع هذا الإشكال أن نضم الدوال \emph{الجزئية} إلى البحث؛ فإننا سنرى
\emph{إمكان} تعداد الدوال الجزئية القابلة للحساب.
\iftag{TMs}{ وقد كاد هذا يتبين لنا من قبل، لأن
آلات تورنغ تقبل تعدادًا منتظمًا.}{} ثم نرجع إلى الحجة القطرية
فنبيّن أين ينقطع الاستدلال إذا أدخلنا الدوال الجزئية.

فما الذي نزيده على الدوال العودية البدائية لنحصل
على الدوال العودية الجزئية كلها؟ لا بد من أمرين:
\begin{enumerate}
\item توسيع تعريف الدوال العودية البدائية حتى يتسع للدوال الجزئية.
\item \emph{إضافة} بند إلى التعريف يُدخل دوال جزئية
جديدة.
\end{enumerate}

أما الأول فيسير: نبدأ بالصفر واللاحق والإسقاطات كما فعلنا، ونأخذ الإغلاق تحت
التركيب والعودية البدائية. وإنما نعدل تعريفي التركيب
والعودية البدائية بحيث يصحان وإن كان بعض الحدود الواردة غير معرّف.
فإذا كانت $f$ و$g$ جزئيتين، كتبنا $f(x)\fdefined$ لنعني أن
$f$ معرفة عند $x$، أي إن $x$ في مجال $f$؛ ونكتب
$f(x)\fundefined$ لنعني العكس، أي إن $f$ غير معرفة عند~$x$.
وأما $f(x)\simeq g(x)$ فمعناه أن $f(x)$ و$g(x)$ إما غير معرفتين
معًا، وإما معرفتان ومتساويتان معًا. ويجري هذا الاصطلاح على الحدود
المركبة أيضًا. ثم نتفق على أنه إذا كانت $h$ و$g_0$،
…،~$g_k$ كلها دوال جزئية، فإن
\[
h(g_0(\vec x),\dots,g_k(\vec x))
\]
تكون معرفة إذا وفقط إذا كانت كل $g_i$ معرفة عند $\vec x$، وكانت $h$
معرفة عند $g_0(\vec x)$، …،~$g_k(\vec x)$. وعلى هذا تبقى صورتا تعريف
التركيب والعودية البدائية كما تقدمتا، ولا يتغير إلا
استبدال «$=$» بـ«$\simeq$».

وأما الإضافة التي تنشئ دوال جزئية جديدة فهي
\emph{مؤثر البحث غير المحدود}. فإن كانت $f(x,\vec z)$ أي دالة جزئية
في الأعداد الطبيعية، عرف $\mu x\;f(x,\vec z)$ بأنه
\begin{quote}
  أصغر $x$ بحيث تكون $f(0,\vec z), f(1,\vec z), \dots, f(x,\vec z)$
  كلها معرفة، وتكون $f(x,\vec z)=0$، إن وجد $x$ كهذا.
\end{quote}
فإن لم يوجد كانت $\mu x\;f(x,\vec z)$ غير معرّفة. وبذلك تتعين
$\mu x\;f(x,\vec z)$ على وجه واحد.

\begin{explain}
وهذا التعريف لا يستند إلى\iftag{TMs}{ آلات تورنغ، ولا إلى}{}
الخوارزميات، ولا إلى أي نموذج حساب معين. غير أن البحث غير المحدود،
كالتركيب والعودية البدائية، له تفسير إجرائي حسابي. فحين نحسب
دالة جزئية، تكون الحجج التي لم تُعرّف عندها هي المداخل
التي لا يقف عليها الحساب. وطريقة حساب $\mu x\;f(x,\vec z)$ هي الآتية:
% تصحيح صياغة كلاسيكية (C224-01): طابق الضمير والخبر الاسم المؤنث «طريقة».
نحسب $f(0,\vec z), f(1,\vec z), f(2,\vec z)$ حتى تعاد قيمة
0. فإن لم يتوقف حساب من الحسابات الوسيطة، لم يتوقف حساب
$\mu x\;f(x,\vec z)$ أيضًا.
\end{explain}

إذا كانت $R(x,\vec z)$ أي علاقة، عرفنا $\mu x\;R(x,\vec z)$ بأنه
$\mu x\;(1\tsub\Char{R}(x,\vec z))$. وبعبارة أخرى، يعيد
$\mu x\;R(x,\vec z)$ أصغر قيمة لـ$x$ بحيث تصدق $R(x,\vec z)$.
فإذا كانت $f(x,\vec z)$ دالة كلية، كان $\mu x\;f(x,\vec z)$ هو نفسه
$\mu x\;(f(x,\vec z)=0)$. أما التعريف الأول فأعم، إذ لا يشترط
أن تكون $f(x,\vec z)$ معرفة في كل موضع (أما الدالة المميزة
لعلاقة فهي كلية دائمًا).

\begin{defn}
مجموعة \emph{الدوال العودية الجزئية} هي أصغر مجموعة من الدوال الجزئية
من الأعداد الطبيعية إلى الأعداد الطبيعية، بمختلف الرتب، تحتوي الصفر
واللاحق والإسقاطات، وتكون منغلقة تحت التركيب والعودية البدائية والبحث
غير المحدود.
\end{defn}

وقد تكون الدالة العودية الجزئية كلية، أعني معرّفة عند كل
حجة؛ فنخص هذه باسم آخر.

\begin{defn}
\ollabel{defn:recursive-fn}
مجموعة \emph{الدوال العودية} هي مجموعة الدوال العودية الجزئية الكلية.
\end{defn}

ويقال أحيانًا «عودية كلية» بدل «عودية»، تنبيهًا إلى أنها معرّفة في كل
موضع.

\end{document}
